\section{Electric correlations and their amplitudes}\label{sec:electric}

The Laplacian estimate of Proposition~\ref{prop:laplacian} determines
the interaction between separated electric charges.  Its integrable
power remainder is essential: it gives a convergent amplitude, rather
than only a logarithmic exponent.  We first obtain that interaction for
small charges, then continue it to the charges that suppress microscopic
loops around the marked points.

We retain the neutral axial correlation
\[
 C(\mathbf t;\mathbf u)
 =\E\exp\!\left(i\sum_{j=1}^p t_jh(u_j)\right),
 \qquad \sum_{j=1}^p t_j=0,
\]
and its potential $H_{\mathbf t}$ from Section~\ref{sec:spectral}.
In the positive winding convention a loop enclosing the set $I$ of
marked points has relative weight
\[
 v_I^+(\mathbf t)=
 \frac{\cos(\lambda+\sum_{j\in I}t_j)}{\cos\lambda}.
\]
The opposite convention has $v_I^-(\mathbf t)=v_I^+(-\mathbf t)$
and gives the same correlation.  In particular,
$C(\mathbf t;\mathbf u)=C(-\mathbf t;\mathbf u)$.
Write
\[
 v(t)=\frac{\cos(\lambda+t)}{\cos\lambda},\qquad
 I_\lambda=(-\pi/2-\lambda,\pi/2-\lambda).
\]
Thus $v(t)>0$ on $I_\lambda$.  As in Section~\ref{sec:radial},
$F_n(v)$ uses the fixed plane cutoff $b_0n$ and fixed lattice rounding.

\subsection{The potential kernel and local charge profiles}

Set $A=1/[2\pi(\pi-\lambda)]$.  The covariance calculation below
identifies this constant in the potential-kernel argument.

\begin{proposition}[Coulomb formula]\label{prop:coulomb}
Fix a number $p$ of marked points.  There are $b,c,\xi>0$ and a
real-analytic function $f$ on a neighborhood of $[-b,b]$, with
$f(0)=0$, such that,
for neutral real charges $\max_j|t_j|\le b$ and distinct points on a
tile axis satisfying
\[
 \max_j|u_j|\le Mn,\qquad
 \min_{i\ne j}|u_i-u_j|\ge n^{1-\xi},
\]
one has
\begin{equation}\label{eq:coulomb}
 \log C(\mathbf t;\mathbf u)
 =\sum_{j=1}^p f(t_j)
   +2A\sum_{i<j}t_it_j\log|u_i-u_j|+O(n^{-c}),
 \qquad f(0)=0.
\end{equation}
Here the same local function $f$ is used for every number of points,
on their common small-charge neighborhood; it is independent of the
geometry and of the point parities.
The constants can depend on $p,M$ and $q$.

For either choice of sign, the corresponding multiplicative statement
has analytic continuation along compact paths from zero in
\[
 \mathcal D_p^\pm
 =\{\mathbf t\in\R^p:\ \sum_jt_j=0,
                         \ \pm t_j\in I_\lambda\text{ for every }j\}.
\]
More precisely, on a sufficiently thin complex neighborhood of each
such compact path within the neutral hyperplane $\sum_jt_j=0$,
there is a holomorphic function $M_p$ such that
\begin{equation}\label{eq:coulomb-continued}
 \frac{C(\mathbf t;\mathbf u)}
 {\displaystyle\prod_{i<j}|u_i-u_j|^{2At_it_j}}
 =M_p(\mathbf t)+O(n^{-c}).
\end{equation}
The positive constants $c$ and $\xi$ may decrease with the path.
Near zero, $M_p(\mathbf t)=\exp(\sum_jf(t_j))$.
The error in \eqref{eq:coulomb-continued} is additive; a relative
error is asserted only where $M_p$ is nonzero.
\end{proposition}

We first derive \eqref{eq:coulomb} for real charges with an unspecified
coefficient in place of $A$; the continuation and the identification of
that coefficient will follow below.  Choose $b$ small enough that
all subset weights in both winding conventions are strictly positive.
The inward coupling of Proposition~\ref{prop:radial}, in its multipoint
form in Proposition~\ref{prop:normalization}, shows that the
unit increments of $H_{\mathbf t}$ near $u_j$ approach a profile
depending only on $t_j$.  Indeed, inside a common surrounding stopping
circuit a marked spin belongs to a loop inside that circuit.  Its
orientation average depends only on whether that loop encloses $u_j$,
and hence only on the single charge $t_j$.  No remote charge enters
this average.  Comparing configurations at separation scales $a$ and
$2a$ gives a power error; summing these errors over successive doublings
constructs a unique limiting increment profile.  Neutrality can always
be maintained by splitting the remote compensating charge into finitely
many small charges.  The same coupling compares two different remote
configurations, including different numbers of compensators: pad them
with zero charges to obtain a common finite number of marked points.
Thus the limit is independent of that choice.

Let $\phi_t(y)$ be the step-two Laplacian of this limiting potential,
with its charge at the origin.  Translation invariance, including an
odd tile translation with the corresponding color exchange, identifies
these profiles for every possible location of the charge.  They also
agree when the remote points are placed on the other tile axis: the
conditional loop law and the orientation average inside the matched
stopping circuit are identical.  This observation will allow us to
compare all four parity classes of the step-two Laplacian.

The profiles have a summable power tail.  To see this at a point $y$,
compare the limiting profile to an experiment whose separation is
$(1+|y|)^B$.  First choose $B$ so large that the inward coupling error
is smaller than $(1+|y|)^{-2-c_2}$ for some $c_2>0$.
Then decrease $b$ so that the factor $n^{Cb}$ in
\eqref{eq:laplacian}, at this value of $n$, does not exhaust the
power gain of that estimate.  It follows that
\begin{equation}\label{eq:profile-tail}
 |\phi_t(y)|\le C(1+|y|)^{-2-c_2}.
\end{equation}
All constants here can be chosen uniformly for the small charges under
consideration, since their positive weights lie in a fixed neighborhood
of one.

Put $g_{\mathbf t}(z)=\Delta_2H_{\mathbf t}(z)$.
The local convergence just established gives the following useful
quantitative form.  After decreasing $b$ and $\xi$, there is $c_3>0$
such that
\begin{equation}\label{eq:profile-approximation}
 \sum_{z\in\Z^2}
 \left|g_{\mathbf t}(z)-
              \sum_j\phi_{t_j}(z-u_j)\right|
 \log(n+|z|+2)=O(n^{-c_3}).
\end{equation}
Here and below the tile coordinate lattice is identified with $\Z^2$.
For completeness, take disjoint balls of radius $n^\varepsilon$ around
the $u_j$.  Within these balls the coupling has a fixed power gain,
with a possible loss $n^{C\xi}$ from the ratio between the configuration
scale and its least separation.  Choose $\xi$ and then $\varepsilon$
small enough that summing over $O(n^{2\varepsilon})$ sites, and
multiplying by a logarithm, still leaves a power gain.  Outside the
balls, \eqref{eq:laplacian} and \eqref{eq:profile-tail} give a bound
of order $n^{Cb-\varepsilon c_1}$, with a possibly smaller exponent
than $c_1$ and an extra logarithm.  Decrease $b$ once more to make this
a negative power.  The same estimates on the unbounded part of the
lattice include the displayed logarithmic weight.  This proves
\eqref{eq:profile-approximation}.

We now apply ordinary lattice potential theory, keeping its parity
issue explicit.  For each $a\in\{0,1\}^2$ let
$\Lambda_a=a+2\Z^2$.  Let $K$ be the potential kernel on $2\Z^2$,
normalized by
\[
 \Delta_2K=\delta_0,\qquad
 K(y)=k\log|y|+k'+O(|y|^{-1})\quad (|y|\longrightarrow\infty),
 \qquad k\ne0.
\]
The kernel normalization, asymptotic and increment estimates follow from
\cite[Proposition~4.4.2, Theorem~4.4.4 and Corollary~4.4.5]{LawlerLimic2010},
after rescaling each parity coset to the integer lattice and adjusting
the generator normalization.  The same kernel is used on each translated
lattice.  By
\eqref{eq:laplacian}, the convolution
\[
 U_a(z)=\sum_{w\in\Lambda_a}K(z-w)g_{\mathbf t}(w),
 \qquad z\in\Lambda_a,
\]
is absolutely convergent, has growth $O_n(1+\log(2+|z|))$, and has
Laplacian $g_{\mathbf t}$.  On the other hand, the exact variogram
bound in Proposition~\ref{prop:spectral} and Cauchy--Schwarz give
\[
 |H_{\mathbf t}(z)-H_{\mathbf t}(z_0)|
 \le |C(\mathbf t;\mathbf u)|^{-1}
       \bigl(\E|h(z)-h(z_0)|^2\bigr)^{1/2}
 =O_n\!\left(1+\sqrt{\log(2+|z-z_0|)}\right).
\]
Thus $H_{\mathbf t}-U_a$ is a sublinearly growing harmonic function
on $\Lambda_a$, and the lattice Liouville theorem makes it constant;
this is the sublinear-growth consequence of the gradient estimate
\cite[Theorem~6.3.8(a)]{LawlerLimic2010}, with the same rescaling.
The tail bound also gives
\[
 U_a(z)=
 \left(\sum_{w\in\Lambda_a}g_{\mathbf t}(w)\right)
                  (k\log|z|+k')+o(1).
\]
The sublogarithmic growth of $H_{\mathbf t}$ therefore forces
\begin{equation}\label{eq:zero-coset-mass}
 \sum_{w\in\Lambda_a}g_{\mathbf t}(w)=0
 \quad\text{for every }a.
\end{equation}
In particular $U_a(z)\to0$ at infinity.  The unit increments of
$H_{\mathbf t}$ tend to zero along the off-axis normals, as proved
in Section~\ref{sec:spectral}.  Evaluating neighboring parity classes
along those normals shows that their four additive constants agree.
Consequently the convolution formula determines all point differences
of $H_{\mathbf t}$, not only differences within one parity class.

Let
\[
 b_a(t)=\sum_{y\in\Lambda_a}\phi_t(y).
\]
Equation~\eqref{eq:profile-approximation}, followed by
\eqref{eq:zero-coset-mass} and $n\to\infty$, implies
\[
 \sum_j b_{a-u_j}(t_j)=0
\]
in every separated neutral experiment; the subscripts are taken modulo
two.  In a two-charge experiment, changing one point by one tile along
the axis shows that $b_a(t)$ is independent of the corresponding
coordinate parity.  Repeating the experiment on the other axis proves
independence of both parities.  Denote the common sum by $b_*(t)$.
Two charges give $b_*(-t)=-b_*(t)$, and three give
$b_*(s+t)=b_*(s)+b_*(t)$ whenever these charges are small.
The uniform bound from \eqref{eq:profile-tail} excludes discontinuous
additive functions.  Hence $b_*(t)=b_1t$ for a constant $b_1$.

Evaluate the convolution formula at $u_i$.  The contribution of its
own profile is
\[
 g_0(t_i)=\sum_{y\in2\Z^2}K(-y)\phi_{t_i}(y),
\]
which depends only on $t_i$.  Every other profile contributes
\[
 b_*(t_j)\bigl(k\log|u_i-u_j|+k'\bigr)
       +O(n^{-c}),
\]
after possibly decreasing $c$.  To justify the power error, split the
profile sum at a small power of $|u_i-u_j|$; use the kernel asymptotic
and its increment bound on the inner part, and
\eqref{eq:profile-tail} on the remaining part.  The error from replacing
$g_{\mathbf t}$ by its profiles is already controlled by
\eqref{eq:profile-approximation}.  Neither this argument nor the
self contribution depends on the point parity.

For $0\le l\le1$ differentiation of the characteristic function gives
the exact identity
\[
 \frac{\dd}{\dd l}\log C(l\mathbf t;\mathbf u)
       =i\sum_i t_iH_{l\mathbf t}(u_i).
\]
The common additive constant disappears by neutrality.  The cross
terms integrate to
$2A\sum_{i<j}t_it_j\log|u_i-u_j|$ with $2A=ikb_1$.
The self terms integrate separately at each point.  The terms involving
$k'$ are also single-point terms, since
$\sum_{i\ne j}t_it_j=-\sum_i t_i^2$.
Absorb them into $f(t_i)$, taking $f(0)=0$.
Uniformity of the preceding estimates for $l\in[0,1]$ proves
\eqref{eq:coulomb} for small real charges.

\subsection{Analytic continuation and the electric constant}

The complex estimates of Section~\ref{sec:radial} turn this local
calculation into a formula at finite charges.  First they determine
the scaling exponent of each positive nesting weight, without yet
determining its amplitude.  On thin complex neighborhoods of compact
subintervals of $I_\lambda$, the functions
\[
 \frac{\log F_n(v(t))}{\log n}
\]
form a locally bounded holomorphic family; logarithms are continued
from real positive weights.  Extract any convergent subsequence.
For small neutral systems, the all-positive comparison in
Proposition~\ref{prop:normalization} bounds the normalized correlation
above and below.  Equation~\eqref{eq:coulomb}, with separations of
order $n$, therefore shows that its limiting exponent $g$ satisfies
\[
 \sum_j g(t_j)=-A\sum_jt_j^2.
\]
The two- and three-charge identities imply
$g(t)=-At^2+Bt$ near zero.  Holomorphic uniqueness extends this identity
along $I_\lambda$.  Since
\[
 v(t)=v(-2\lambda-t),
\]
the same symmetry holds for $g$, and it forces $B=-2\lambda A$.
Every subsequential exponent is thus the same.  We have proved,
locally uniformly also in those complex neighborhoods,
\begin{equation}\label{eq:nest-log-exponent}
 F_n(v(t))=\exp\{\beta(t)\log n+o(\log n)\},
 \qquad \beta(t)=-A(t^2+2\lambda t).
\end{equation}

Divide $C(\mathbf t;\mathbf u)$ by the cross-distance product in
\eqref{eq:coulomb-continued}.  Neutrality gives
$\sum_j\beta(t_j)=-A\sum_jt_j^2$, so
Proposition~\ref{prop:normalization} and
\eqref{eq:nest-log-exponent} bound this quotient by $n^{o(1)}$
on compact complex neighborhoods with macroscopic separations.
If the least separation is only $n^{1-\xi}$, the bound is
$n^{o(1)+C\xi}$.  These are bounds for the quotient itself; no
nonvanishing claim about the correlation is needed here.

Here is the quantitative continuation argument.  Fix a reference
configuration with macroscopic separations and fixed parities.  Compare
its quotients at sizes $n$ and $N$, uniformly for $n\le N\le2n$.
Equation~\eqref{eq:coulomb} makes their difference $O(n^{-c})$ on
a small real neighborhood of zero.  Apply the two-constants theorem
in disks slit along a compact real segment, successively in the
$p-1$ independent charge variables.  A negative power on the segment
and an $n^{o(1)}$ bound on the other boundary give a negative power
on a smaller complex neighborhood.  A finite chain of overlapping
disks carries this estimate along any prescribed compact path in
$\mathcal D_p^+$.  The resulting exponent may be smaller, but remains
positive.  Summing the comparison over dyadic scales, and using its
uniformity for every $N\in[n,2n]$, gives a holomorphic limit $M_p$
through all integers with a power error.

Next subtract the reference quotient from the quotient of any allowed
configuration.  Their difference has the same small-charge power bound
by \eqref{eq:coulomb}, while its bound along the continuation path is
$n^{o(1)+C\xi}$.  Choose $\xi$ sufficiently small after fixing that
path.  The same finite sequence of two-constants estimates still leaves
a negative power.  This proves the geometry-uniform statement in
\eqref{eq:coulomb-continued}.  The opposite winding convention proves
the assertion on $\mathcal D_p^-$.  The initial amplitude is therefore
real analytic.  To obtain the same conclusion for $f$, compare the
logarithmic amplitudes of the neutral triples $(s,t,-s-t)$ and pairs
$(s+t,-s-t)$.  Their difference is the analytic function
$f(s)+f(t)-f(s+t)$.  Averaging this identity in $s$ over a fixed small
interval first shows that the continuous function $f$ is continuously
differentiable; differentiation at $s=0$ then expresses $f'(t)$ as a
constant minus an analytic function.  Hence $f$ is real analytic.

It remains to identify $A$.  Take four distinct collinear continuum
points $a,b,c,d$ and the neutral charges $(-s,s,-t,t)$ at their lattice
approximations.  The mixed derivative at zero of $\log C$ is minus
the covariance of the two height increments.  In
\eqref{eq:coulomb} it kills all single-point terms and yields
\[
 -2A\log\frac{|a-c|\,|b-d|}{|a-d|\,|b-c|}
\]
for that covariance in the limit.  Differentiation of the limit is
legitimate by the locally uniform analytic convergence just proved.
Input~\ref{inp:covariance}, with
$G(z,w)=-(2\pi)^{-1}\log|z-w|$ and
$c=2\cos(\lambda/2)$ in the six-vertex parametrization, gives
\begin{equation}\label{eq:stiffness}
 2A=\frac{\sigma(2\cos(\lambda/2))^2}{2\pi}
    =\frac{1}{\pi(\pi-\lambda)},
 \qquad A=\frac{1}{2\pi(\pi-\lambda)}.
\end{equation}
This is the separated-endpoint covariance assertion of
\cite[Theorem~1.1, Part~(iii)]{GFFCompanion}; at the present parameters
it lies in the established interval $\sqrt3\le c<2$ of
\cite[Theorem~2.8]{DKLM}.  Only the covariance statement at that scope
is used here.  This completes the proof of
Proposition~\ref{prop:coulomb}.

\subsection{Pure nesting amplitudes and nonzero special charges}

\begin{proposition}[Individual nesting amplitudes]\label{prop:pure-nest}
There is $\varepsilon_0>0$ and a nonvanishing holomorphic function
$a(t)$ near $[-\varepsilon_0,\varepsilon_0]$, positive for real $t$,
such that
\[
 F_n(v(t))=a(t)n^{\beta(t)}(1+o(1))
\]
locally uniformly on a complex neighborhood of that interval.
The limit uses the fixed cutoff and rounding in the definition of
$F_n$.  No rate is asserted for this individual amplitude.
\end{proposition}

\begin{proof}
Set
\[
 L_n(t)=\log F_n(v(t))-\beta(t)\log n,\qquad L_n(0)=0.
\]
For fixed macroscopic two- and three-point configurations,
Proposition~\ref{prop:normalization} gives a holomorphic limit of
$C/\prod_jF_n(v(t_j))$, nonzero near zero because its value at zero
is one.  Proposition~\ref{prop:coulomb} gives the corresponding
nonzero limit after division by $n^{\sum_j\beta(t_j)}$.
Taking the logarithms normalized at zero shows that
$\sum_jL_n(t_j)$ converges holomorphically for each such neutral
system.  Subtract the identities for $(s,t,-s-t)$ and $(s+t,-s-t)$.
It follows that
\[
 L_n(s)+L_n(t)-L_n(s+t)
\]
converges near $(0,0)$.  Differentiation at $s=0$ shows that
$L_n'(t)-L_n'(0)$ converges.  Hence, with the real number
$b_n=L_n'(0)$, the functions $L_n(t)-b_nt$ converge holomorphically
near zero.  Neutral correlations alone have left precisely this
linear ambiguity.

To determine $b_n$, choose one charge $t$ in a small interval about
$-2\lambda$ and choose sufficiently many other charges
$s_1,\ldots,s_k$ in the neighborhood just considered, with
$\sum_js_j=-t$.  They can be chosen strictly inside that neighborhood.
All these individual charges belong to $I_\lambda$, and the entire
neutral parameter set connects to zero inside $\mathcal D_{k+1}^+$.
For this fixed macroscopic configuration both the continued amplitude
in Proposition~\ref{prop:coulomb} and the limit in
Proposition~\ref{prop:normalization} are analytic and are not
identically zero.  Their restrictions cannot vanish on an open real
parameter set.  Choose real parameters in the indicated open set where
both are nonzero.  Taking absolute values and ordinary real logarithms
then shows that
\[
 L_n(t)+\sum_jL_n(s_j)
\]
converges.  Put $t'=-2\lambda-t$, which is near zero.
The identities $v(t')=v(t)$ and $\beta(t')=\beta(t)$ give
$L_n(t)=L_n(t')$ exactly.  Substituting the already convergent
remainders $L_n(z)-b_nz$ leaves
\[
 b_n\left(t'+\sum_js_j\right)=b_n(t'-t).
\]
The fixed coefficient $t'-t$ is nonzero, so $b_n$ converges.
Exponentiating the resulting locally uniform limit of $L_n$ proves
the proposition.
\end{proof}

We next specify the charges used to detect microscopic arms.  A
single group consists of one charge $\alpha=\pi/2-\lambda$ and
at least two small negative charges whose sum is $-\alpha$.
Choose sufficiently many compensators that each lies, with its
negative, in the interval of Proposition~\ref{prop:pure-nest}.
The number of charges is now fixed.  They need not all lie in the
small-system neighborhood used in the proof of
\eqref{eq:laplacian} for that number of points.

\begin{proposition}[Two special core charges]\label{prop:special-charges}
Two separately neutral groups of the preceding kind can be chosen
so that their combined amplitude in
\eqref{eq:coulomb-continued} is nonzero.  The formula then has a
relative power error, including when the two $\alpha$-points have
opposite colors and least separation $r\ge n^{1-\xi}$, for a
sufficiently small fixed $\xi>0$.

Place the first core at the origin, the second at odd axial separation
$r\asymp n^{1-\xi}$, and all compensators at distinct positions
$n z_j+O(1)$, with fixed nonzero axial $z_j$.  Write
$B_* =\sum_{j\,\mathrm{remote}}\beta(t_j)$.
There is a constant $M_*\ne0$, depending only on those charges,
their macroscopic positions, and the fixed microscopic conventions,
such that
\begin{equation}\label{eq:special-close-pair}
 C_r(n)=M_*n^{B_*-x}r^{2A\alpha^2}(1+O(n^{-c}))
       =M_*n^{B_*}r^{-2x_1}(r/n)^x(1+O(n^{-c})),
 \qquad x_1=-\beta(\alpha).
\end{equation}
The same conclusion is uniform for
$n^{1-\xi}\le r\le n^{1-\xi'}$, with fixed $0<\xi'\le\xi$,
after decreasing $c$.  For $r=an$ with fixed small $a>0$, the full
Coulomb product applies; its remaining geometric factor tends to
$M_*$ as $a\downarrow0$ and stays bounded above and away from zero.
\end{proposition}

\begin{proof}
For one group use the opposite winding convention.  Every subset
charge lies between $-\alpha$ and $\alpha$, so all scalar weights
are nonnegative.  The only possible zero is a loop enclosing every
compensator and excluding the $\alpha$-point.  Every singleton weight
is strictly positive: at the core this uses $\lambda>0$, and at
the compensators it follows from their smallness.  The group thus
lies strictly inside $\mathcal D_p^-$, where the continued formula
has a power error.

Its real correlation compares above and below to the product of its
opposite-convention singleton factors $F_n(v(-t_j))$.
For the lower comparison, impose surrounding primal circuits in the
separated point patches and connect them by primal paths outside the
smaller cores.  At least two compensators are used.  A loop enclosing
all compensators but not the core could not lie within one patch and
could not cross this connected collection of circuits and paths;
therefore the zero-weight loops are excluded.  The circuit and
connection event costs a fixed positive probability by RSW and
collar comparison.  More explicitly, first integrate the positive
singleton tilts in the smaller cores, leave full collars to the
connecting event, and then use the passage estimates and the local
tilt bounds of Section~\ref{sec:radial}.  All remaining charged
macroscopic loop counts have fixed exponential moments under this
proxy law.  Truncate them at a sufficiently large fixed number, losing
less than half the lower probability of the connecting event.
On the retained event every nonzero scalar factor is bounded below,
which proves the lower comparison.  The upper comparison follows
from the same passage moments.

By \eqref{eq:nest-log-exponent} this comparison fixes the logarithmic
exponent of the group correlation.  After division by its Coulomb
product the result is bounded below by $n^{-o(1)}$.
If its amplitude vanished, \eqref{eq:coulomb-continued} would instead
bound that quotient by $O(n^{-c})$, a contradiction.  Thus the
single-group amplitude is nonzero.  For two independently neutral
groups, the identity
\[
 M_{p+p'}(\mathbf t,\mathbf s)=M_p(\mathbf t)M_{p'}(\mathbf s)
\]
holds near zero by the single-point form of
\eqref{eq:coulomb}, and hence throughout the continuation with
each group separately neutral.  The combined amplitude is nonzero
as well.  The combined positive-convention scalar weights may have
either sign; nonvanishing of the combined amplitude has been proved
using the opposite convention and analytic continuation.

Finally separate the core-core factor $r^{2A\alpha^2}$ in the
Coulomb product.  Every other distance divided by $n$ is a fixed
nonzero number plus $O(r/n+n^{-1})$.  Neutrality and
\eqref{eq:stiffness} give
\[
 2x_1+2A\alpha^2
 =4A\alpha(\alpha+\lambda)
 =\frac{\pi-2\lambda}{2(\pi-\lambda)}=x.
\]
The exponent of $n$ in the remaining product is consequently
$B_*-x$.  When $r/n$ tends to zero at a power rate its geometric
factor has a power-rate limit $M_*$, proving
\eqref{eq:special-close-pair}.  Keeping $r/n=a$ instead gives the
last assertion directly from the uncontracted Coulomb product.
\end{proof}

\subsection{Arm estimates needed for microscopic comparison}

Let $\pi_1(R)$ be the plane probability that a point is connected in
its own color to distance $R$, with a fixed bounded microscopic inner
scale.  Let $\pi_2(r,R)$ be the plane probability of a primal arm and
a dual arm across the concentric square annulus with radii $r,R$.
All connections here are actual lattice paths.  Bounded changes of
these radii are harmless by the arm estimates in Input~\ref{inp:fk}.

\begin{lemma}[Sharp one-arm comparison]\label{lem:one-arm}
There are constants $0<c<C<\infty$ such that
\begin{equation}\label{eq:one-arm}
 cR^{-x_1}\le\pi_1(R)\le CR^{-x_1},
 \qquad x_1=-\beta(\alpha),
\end{equation}
for all radii above the fixed microscopic scale.
\end{lemma}

\begin{proof}
Use one of the neutral single-core groups above at mutual separations
of order $R$, now in the positive winding convention.  Its singleton
core weight is $v(\alpha)=0$.  Thus no loop wholly inside its core
patch may enclose the point.  In the switched-region description this
is exactly connection of the point's color to the patch rim, up to a
bounded resizing of a whole-tile box.  All weights of this single group
are nonnegative.  Localize the small positive compensator weights in
their disjoint patches.  Passage bounds through collars give the upper
comparison
\[
 C(\mathbf t;\mathbf u)
 \le C\pi_1(R)\prod_{j\,\mathrm{remote}}F_R(v(t_j)).
\]

For the reverse comparison, give the core arm a separated landing and
continue it through a fixed corridor to a circuit of the same color
around one compensator patch.  Any loop enclosing the core point must
then enclose that circuit and its compensator as well: otherwise it
would cross the connected color path.  Hence no remaining loop has
the zero core-only weight.  The landing extension, corridor, and
circuit cost a bounded factor by Input~\ref{inp:fk}, uniformly after
the positive tilts inside the compensator cores, since full collars
separate those tilts from the construction.  The macroscopic loop
counts still have every fixed exponential moment with a bounded cost
relative to the arm probability.  To apply that bound, first cut the
arm requirement back through a fixed-factor buffer, use arm
quasimultiplicativity and localization, and leave space between its
support and the passage buffers.  The local tilt bounds deal with the
compensator cutoffs.  A fixed truncation of these counts then retains
a positive fraction of the constructed event and bounds all nonzero
weights below.  Therefore
\[
 C(\mathbf t;\mathbf u)
 \asymp\pi_1(R)\prod_{j\,\mathrm{remote}}F_R(v(t_j)).
\]
The nonzero Coulomb amplitude makes the left side comparable to
$R^{\sum_j\beta(t_j)}$, and
Proposition~\ref{prop:pure-nest} supplies the compensator powers.
Canceling them proves \eqref{eq:one-arm}.
\end{proof}

\begin{lemma}[Logarithmic two-arm exponent]\label{lem:two-arm-log}
For every $\delta>0$ there is $C_\delta<\infty$ such that
\begin{equation}\label{eq:two-arm-log}
 C_\delta^{-1}(r/R)^{x+\delta}
 \le\pi_2(r,R)\le C_\delta(r/R)^{x-\delta},
 \qquad R\ge r\ge r_0,
\end{equation}
where $r_0$ is a fixed microscopic scale.
\end{lemma}

\begin{proof}
In a Dobrushin square of size $R$, place the test disk at its center.
For $r\le R/C$, the probability that the distinguished interface enters
this disk compares to $\pi_2(r,R)$, allowing fixed factors in the
radii.  A disk visit requires both color paths to distance comparable
to $R$, giving the upper comparison through full collars.  For the
lower comparison, separate both ends of the annular arms by the
landing-extension result in Input~\ref{inp:fk}.  Continue the outer
primal arm to a flat portion of the wired wall and the outer dual arm
to a flat portion of the complementary dual wall, away from the marked
corners, through disjoint corridors.  Up-to-wall crossings cost a
bounded factor.  Continue both inner ends into a fixed enlargement of
the disk.  The distinguished strand separating these two boundary
regions must then enter that enlargement.  These are the usual
landing and localization arguments of \cite[Section~6]{DMT}; the
paths reach the actual wall neighborhoods and do not use boundary
identification as an annular crossing.  Removing the wired perimeter
edges from the sampled set does not alter the argument, since their
endpoints are already identified and those edges are declared open
when drawing the strand.

Fix a large annulus ratio $K$ and let the inner radius tend to infinity.
Input~\ref{inp:trace} applies to these square approximations.
Open and closed disk-hitting events, with a fixed enlargement or
contraction of the radius, give upper and lower limiting bounds even
without selecting a radius of continuity.  The Euclidean SLE Green
function estimate \cite[Theorem~2.3, Equation~(12)]{LawlerRezaeiBasic}, applied at
the fixed interior center of the square by local conformal distortion,
gives a positive asymptotic proportional to $K^{-x}$ as
$K\to\infty$.  Its hypothesis $\kappa<8$ holds here.  The constants
in the preceding arm-to-disk comparisons do not depend on $K$.
Consequently there are constants $c,C>0$, independent of large $K$,
such that, for each fixed $K$ and all sufficiently large $r$,
\[
 cK^{-x}\le\pi_2(r,Kr)\le CK^{-x}.
\]

Let $Q$ be a quasimultiplicativity constant.  Choose $K$ so large
that the fixed factors $Q$, $c^{-1}$, and $C$ in a product over one
annulus are absorbed by $K^{\delta/2}$.  Iterate the last display
on annuli of ratio $K$, once their inner radii exceed the corresponding
threshold.  Quasimultiplicativity gives the powers $x\pm\delta/2$.
A final annulus of bounded ratio is handled by RSW, and the finitely
many scales below the threshold are absorbed in $C_\delta$ using
finite energy and fixed-factor arm extensions.  Enlarging the constant
and relaxing to $x\pm\delta$ proves \eqref{eq:two-arm-log} for all
the stated radii.  This argument supplies the logarithmic exponent;
the sharp two-arm comparison will be obtained from the microscopic
calibration in Section~\ref{sec:calibration}.
\end{proof}
